% !TeX root = ../main.tex
\section{Badania symulacyjne}
\label{sec:simulation_results}

% ------------------------------------------------------------------------
% Rozdział 4: Badania symulacyjne efektywności detekcji radaru
% 4.1 Scenariusze testowe (parametryzacja, stała moc AP 100 mW EIRP, degradacja szumowa 9 dB dla 160 MHz)
% 4.2 Rozdzielczość odległościowa a szerokość pasma (E1: 20, 40, 80, 160 MHz)
% 4.3 Długość symbolu OFDM a prędkość jednoznaczna (E2: 802.11a vs 802.11ax, aliasing Dopplera)
% 4.4 Wpływ rzędu modulacji MCS i poziomu szumu na estymację ZF (E3: wzmocnienie szumu zeta_ZF, kontrast PNFR)
% 4.5 Tłumienie stacjonarnego clutteru (E4: modele TGax A, C, E, filtr MTI w czasie wolnym)
% 4.6 Skuteczność algorytmu Coherent CLEAN (E5: duża dynamika mocy >30 dB, odsłanianie celów słabych, adaptacyjny próg MAD, Monte Carlo)
% ------------------------------------------------------------------------

\subsection{Scenariusze testowe}
\label{subsec:sim_scenarios}

\subsection{Rozdzielczość odległościowa a szerokość pasma}
\label{subsec:sim_bandwidth_resolution}

\subsection{Długość symbolu OFDM a prędkość jednoznaczna}
\label{subsec:sim_doppler_aliasing}

\subsection{Wpływ rzędu modulacji MCS i poziomu szumu na estymację ZF}
\label{subsec:sim_mcs_noise_enhancement}

\subsection{Tłumienie stacjonarnego clutteru}
\label{subsec:sim_clutter_suppression}

\subsection{Skuteczność algorytmu Coherent CLEAN}
\label{subsec:sim_clean_performance}
